% 来源：《理论力学》课堂笔记手写扫描件 第61页全页、第62页上半部分（第九章）
\section{质点动力学的基本方程}

动力学研究物体的机械运动与作用力之间的关系。

\subsection{动力学基本定律}

\parahead{Ⅰ. 牛顿第一定律（惯性定律）}

不受力作用的质点，将保持静止或做匀速直线运动状态，即：物体具有惯性。

\parahead{Ⅱ. 牛顿第二定律（力与加速度关系定律）}

% 待核对：原稿作「即得明得摩尔著名的 F=ma」，中间二字字形难辨，疑为「牛顿」，此处取「牛顿」。
牛顿第二定律可表述为
\begin{equation}\label{eq:10-newton-second}
  \frac{d}{dt}\left(m\vec{v}\right)=\vec{F},
\end{equation}
即得牛顿著名的 $\vec{F}=m\vec{a}$。
% 待核对：原稿作「成立的光条件」，疑为「前提条件／充分条件」，此处取「前提条件」。
该式成立的前提条件是在经典力学框架下物体质量不变。

\parahead{Ⅲ. 牛顿第三定律（作用与反作用定律）}

即静力学公理四。它不仅适用于平衡的物体，也适用于运动的物体，且与参考系的选取无关。

\subsection{质点的运动微分方程}

由牛顿第二定律，质点的运动微分方程（矢量形式）为
\begin{equation}\label{eq:10-motion-ode}
  m\frac{d^2\vec{r}}{dt^2}=\sum\vec{F}.
\end{equation}

\subsection{在直角坐标轴上的投影}

设矢径 $\vec{r}$ 在直角坐标轴上的投影分别为 $x,y,z$，力 $\vec{F}$ 在坐标轴上的投影分别为 $F_x,F_y,F_z$，则运动微分方程的投影式为
\begin{equation}\label{eq:10-projection-xyz}
  m\frac{d^2x}{dt^2}=\sum F_x,\qquad
  m\frac{d^2y}{dt^2}=\sum F_y,\qquad
  m\frac{d^2z}{dt^2}=\sum F_z.
\end{equation}

\subsection{在自然轴上的投影}

已知加速度沿自然轴的分解
\begin{equation}\label{eq:10-accel-natural}
  \vec{a}=a_t\vec{e}_t+a_n\vec{e}_n,\qquad a_b=0,
\end{equation}
则在自然轴上有投影式
\begin{equation}\label{eq:10-projection-natural}
  m\frac{dv}{dt}=\sum F_t,\qquad
  m\frac{v^2}{\rho}=\sum F_n,\qquad
  0=\sum F_b.
\end{equation}
